\documentclass[11pt]{article}
\usepackage[margin=1in]{geometry}
\usepackage{amsmath,amssymb,amsfonts,mathtools}
\usepackage{booktabs}
\usepackage{xcolor}
\usepackage{hyperref}
\usepackage{enumitem}
\usepackage{array}
\usepackage{tabularx}

\hypersetup{colorlinks=true,linkcolor=blue,citecolor=blue,urlcolor=blue}

\newcommand{\R}{\mathbb{R}}
\newcommand{\1}{\mathbf{1}}
\newcommand{\softmax}{\operatorname{softmax}}
\newcommand{\softplus}{\operatorname{softplus}}

\title{\textbf{TP2-MV-ALiBi Transformer}\\[4pt]
\large Architecture Explanation with Formulas}
\author{}
\date{}

\begin{document}
\maketitle
\tableofcontents
\newpage

%%────────────────────────────────────────────────────────────────
\section{What Problem the Transformer Solves}

For each TP2 episode on a given day, we want to answer:

\begin{center}
\textit{``Will trading strategy T1 (or K2) on this violation make money by expiry?''}
\end{center}

A standard model (MLP, LightGBM) looks at the current episode's 48 features and
gives a probability. The transformer does something extra: it asks

\begin{center}
\textit{``What happened historically in episodes that looked similar to this one?''}
\end{center}

and uses that historical context to improve the prediction.

The key structural insight used throughout: both viable strategies T1 and K2
\textbf{short the shorter-dated D-leg}. T2 and K1 short the longer-dated C-leg
and are structural losers. The strategy universe is therefore restricted to
$r \in \{T1, K2\}$ as a theoretical prior.

%%────────────────────────────────────────────────────────────────
\section{Token Construction}

Each (episode $i$, strategy $r$) pair becomes one \textbf{token} --- a vector
of numbers the transformer reads.

\[
  u_{i,r}
  =
  \bigl[\,
    x_{i,r};\;
    \mathrm{emb}(r);\;
    \mathrm{emb}(\mathrm{month}_i);\;
    \mathrm{emb}(\mathrm{dayofweek}_i)
  \,\bigr]
\]

where:
\begin{itemize}[leftmargin=*]
  \item $x_{i,r} \in \R^{48}$ --- the 48 features selected from $\mathcal{F}^\star$
        (option surface, short-D-leg risk, liquidity, lifecycle pressure).
  \item $\mathrm{emb}(r)$ --- a learned embedding for whether the strategy is T1 or K2.
  \item Month and day-of-week --- seasonality signals.
  \item \texttt{signal\_year} is \textbf{explicitly excluded} to prevent fold leakage.
\end{itemize}

The token is projected to a hidden dimension $d = 64$:
\[
  h_{i,r}^{(0)}
  =
  \mathrm{MLP}_{\mathrm{emb}}\!\bigl(\mathrm{LayerNorm}(u_{i,r})\bigr)
  \;\in\; \R^{64}.
\]

The MLP embedding uses the structure:
\[
  \mathrm{Linear}(48 \to 64) \;\to\; \mathrm{LayerNorm} \;\to\; \mathrm{GELU}
  \;\to\; \mathrm{Dropout}(0.10) \;\to\; \mathrm{Linear}(64 \to 64).
\]

%%────────────────────────────────────────────────────────────────
\section{Standard Self-Attention (Baseline)}

In standard attention with $H = 4$ heads, each of dimension
$d_h = 64 / 4 = 16$, the attention score between episode $i$ (query)
and historical episode $j$ (key) in head $h$ is:

\[
  a_{ij}^{(h)}
  =
  \frac{(Q_i^{(h)})^\top K_j^{(h)}}{\sqrt{d_h}}
\]

Attention weights:
\[
  \alpha_{ij}^{(h)}
  =
  \frac{\exp(a_{ij}^{(h)})}{\displaystyle\sum_{k} \exp(a_{ik}^{(h)})}
\]

Context vector (weighted sum of value vectors):
\[
  \tilde{h}_i^{(h)}
  =
  \sum_j \alpha_{ij}^{(h)} V_j^{(h)}
\]

\textbf{Problem with standard attention for TP2:} All historical episodes
are treated as equally relevant regardless of when they happened, what the
option surface looked like, or how liquid the market was. A violation from
3 years ago in a completely different volatility regime gets the same
consideration as one from last week in an identical setup.

%%────────────────────────────────────────────────────────────────
\section{Original ALiBi --- Scalar Distance Penalty}

ALiBi (Press, Smith \& Lewis, ICLR 2022) adds a \textbf{penalty proportional
to time distance} to the attention score, replacing positional embeddings:

\[
  a_{ij}^{(h)}
  =
  \frac{(Q_i^{(h)})^\top K_j^{(h)}}{\sqrt{d_h}}
  - m_h \cdot (t_i - t_j)
\]

where $m_h > 0$ is a fixed head-specific slope. Larger slope = stronger
preference for recent episodes. The penalty grows linearly with how far back
you look.

\textbf{Limitation for TP2:} Scalar time distance alone does not capture what
actually matters. A violation from 5 years ago with an identical option surface
configuration and the same short-D-leg moneyness is more relevant than a
violation from last month with a completely different setup.

%%────────────────────────────────────────────────────────────────
\section{TP2-MV-ALiBi --- Multivariate Distance Penalty}

\subsection{Core Idea}

We replace the scalar time distance with a \textbf{vector of 5 TP2-domain
distances}:

\[
  d_{ij}
  =
  \bigl[\,
    D_{\mathrm{time}}(i,j),\;
    D_{\mathrm{surface}}(i,j),\;
    D_{\mathrm{short\text{-}D}}(i,j),\;
    D_{\mathrm{liq}}(i,j),\;
    D_{\mathrm{life}}(i,j)
  \,\bigr]^\top
  \;\in\; \R^5
\]

Each attention head gets a \textbf{slope vector} $m_h \in \R^5$:

\[
  b_{ij}^{(h)} = -m_h^\top d_{ij}
\]

\subsection{Full Attention Logit}

\[
  \boxed{
  a_{ij}^{(h)}
  =
  \underbrace{\frac{(Q_i^{(h)})^\top K_j^{(h)}}{\sqrt{d_h}}}_{\text{feature similarity}}
  -\;
  \underbrace{m_h^\top d_{ij}}_{\text{MV-ALiBi penalty}}
  +\;
  \underbrace{\rho_h^{\mathrm{same}} \cdot \1\{r_i = r_j\}}_{\text{same-strategy bonus}}
  +\;
  \underbrace{\rho_h^{\mathrm{cross}} \cdot \1\{r_i \neq r_j\}}_{\text{cross-strategy bonus}}
  }
\]

where:
\begin{itemize}[leftmargin=*]
  \item $\rho_h^{\mathrm{same}}$ --- learned bonus when both tokens are the same
        strategy (T1$\leftrightarrow$T1 or K2$\leftrightarrow$K2).
  \item $\rho_h^{\mathrm{cross}}$ --- learned bonus for cross-strategy comparison
        (T1$\leftrightarrow$K2).
\end{itemize}

\subsection{Slopes are Guaranteed Positive (Softplus Parameterization)}

The slopes $m_h$ are parameterized as:
\[
  m_h = \softplus(\tilde{m}_h),
  \qquad
  \softplus(x) = \log(1 + e^x)
\]
where $\tilde{m}_h$ is the learned parameter. This guarantees $m_h \geq 0$,
which preserves the ALiBi property: \textbf{larger distances always reduce
attention weight}.

\subsection{Attention Weights and Context}

\[
  \alpha_{ij}^{(h)}
  =
  \frac{\exp(a_{ij}^{(h)})}{\sum_{k < i} \exp(a_{ik}^{(h)})}
\]

Padding positions (empty history) are masked with $-10^9$ before softmax.

\[
  c_i
  =
  \sum_{j < i} \alpha_{ij} V_j
  \;\in\; \R^{64}
\]

%%────────────────────────────────────────────────────────────────
\section{The Five Distance Functions}

\subsection{D1 --- Time Distance}
How old is the historical episode?
\[
  D_{\mathrm{time}}(i,j)
  =
  \log(1 + \Delta t_{ij}),
  \qquad
  \Delta t_{ij} = t_i - t_j \;\text{(in days)}
\]
Log-scaled so the penalty does not explode for very old episodes.

\subsection{D2 --- Option Surface Distance}
Where on the moneyness--maturity surface does the violation occur?
\[
  D_{\mathrm{surface}}(i,j)
  =
  \bigl\|W_s(s_i - s_j)\bigr\|_1
\]
\[
  s_i
  =
  \bigl[
    \log(K_A/S_i),\;
    \log(K_B/S_i),\;
    \log(K_C/S_i),\;
    \log(K_D/S_i),\;
    \log(1 + \mathrm{DTE}_{T1,i}),\;
    \log(1 + \mathrm{DTE}_{T2,i})
  \bigr]
\]

\subsection{D3 --- Short-D-Leg Distance}
Risk profile of the leg that \textbf{both T1 and K2 short}:
\[
  D_{\mathrm{short\text{-}D}}(i,j)
  =
  \bigl\|W_D(z_{D,i} - z_{D,j})\bigr\|_1
\]
\[
  z_{D,i}
  =
  \bigl[
    \log(K_D/S_i),\;
    \Delta_{D,i},\;
    \Gamma_{D,i},\;
    \Theta_{D,i},\;
    \sigma_{D,i},\;
    \mathrm{DTE}_{D,i},\;
    \mathrm{spread}_{D,i}
  \bigr]
\]
This distance is the structural heart of the model: T1 and K2 are both
profitable because they short the same short-dated D-leg. Historically similar
short-D-leg configurations should predict similar outcomes.

\subsection{D4 --- Liquidity Distance}
How tradeable is the episode?
\[
  D_{\mathrm{liq}}(i,j)
  =
  \bigl\|W_\ell(\ell_{i,r} - \ell_{j,r})\bigr\|_1
\]
\[
  \ell_{i,r}
  =
  \bigl[
    \mathrm{max\ spread\ pct}_{i,r},\;
    \mathrm{mean\ spread\ pct}_{i,r},\;
    \mathrm{min\ volume}_{i,r},\;
    \mathrm{roundtrip\ cost}_{i,r}
  \bigr]
\]

\subsection{D5 --- Lifecycle Distance}
Trajectory and persistence of the TP2 violation before signal date:
\[
  D_{\mathrm{life}}(i,j)
  =
  \bigl\|W_c(c_i - c_j)\bigr\|_1
\]
\[
  c_i
  =
  \bigl[
    \mathrm{lc\_pressure}_i,\;
    \mathrm{lc\_pressure\_slope}_i,\;
    \mathrm{lc\_observable\_share}_i
  \bigr]
\]

In the implementation, $W_s, W_D, W_\ell, W_c$ are computed as
\textbf{mean absolute difference} across that feature block --- simple,
interpretable, and parameter-free.

%%────────────────────────────────────────────────────────────────
\section{Prediction Heads}

Combined representation after attention:
\[
  z_i = \bigl[h_i^{(0)};\; c_i\bigr] \;\in\; \R^{128}
\]

\subsection{Binary Classification Head (Main --- Used for Trading)}
\[
  \hat{p}_{i,r}
  =
  \sigma\!\left(
    \mathrm{MLP}_{\mathrm{cls}}\!\left(
      \mathrm{LayerNorm}(z_i)
    \right)
  \right)
\]
$\sigma$ is the sigmoid function. $\hat{p}_{i,r}$ is the probability that
strategy $r$ on episode $i$ is profitable at expiry.

\subsection{Return Auxiliary Head (Diagnostics Only)}
\[
  \hat{\mu}_{i,r}
  =
  \mathrm{MLP}_{\mathrm{ret}}\!\left(
    \mathrm{LayerNorm}(z_i)
  \right)
\]
Predicts the normalized realized return. Used for diagnostics and error
analysis. \textbf{Not used for trading decisions.}

\subsection{Trading Decision Rule}
\[
  r_i^\star
  =
  \arg\max_{r \in \{T1, K2\}} \hat{p}_{i,r},
  \qquad
  \text{trade only if } \max_{r} \hat{p}_{i,r} > \theta
\]
The threshold $\theta$ is tuned on the validation year only. Test year is
evaluated once.

%%────────────────────────────────────────────────────────────────
\section{Training Loss}

\[
  \mathcal{L}
  =
  \mathcal{L}_{\mathrm{BCE}}
  +
  \lambda_R \cdot \mathcal{L}_{\mathrm{Huber}}
\]

\subsection{Binary Cross-Entropy (Main Loss)}
With class-weight $w^+ = (1-\bar{y})/\bar{y}$ to handle class imbalance
($\bar{y}$ is the training positive rate):

\[
  \mathcal{L}_{\mathrm{BCE}}
  =
  -\sum_i
  \Bigl[
    w^+ y_i \log \hat{p}_i
    +
    (1 - y_i) \log(1 - \hat{p}_i)
  \Bigr]
\]

\subsection{Huber Return Loss (Auxiliary, $\lambda_R = 0.02$)}
\[
  \mathcal{L}_{\mathrm{Huber}}
  =
  \sum_i \mathrm{Huber}_{0.5}\!\left(R_i^{\mathrm{expiry}} - \hat{\mu}_i\right)
\]

\[
  \mathrm{Huber}_\delta(r)
  =
  \begin{cases}
    \tfrac{1}{2}r^2 & |r| \leq \delta \\
    \delta\bigl(|r| - \tfrac{1}{2}\delta\bigr) & |r| > \delta
  \end{cases}
\]

The Huber loss is less sensitive to the large loss outliers that characterize
failing T2/K1-style episodes (which are excluded by the mechanical gate but
can occur in borderline cases). $\lambda_R = 0.02$ keeps it strictly
auxiliary.

\textbf{Pairwise ranking loss is intentionally excluded.} Cross-episode
ranking is not conceptually clean for this trading problem: two episodes on
different days with different market conditions are not comparable by a
ranking objective.

%%────────────────────────────────────────────────────────────────
\section{Isotonic Probability Calibration}

After training, raw probabilities are passed through an isotonic regression
fitted on the validation year:

\[
  \hat{p}^{\mathrm{cal}} = \mathrm{IsotonicRegression}(\hat{p}^{\mathrm{raw}})
\]

\begin{itemize}[leftmargin=*]
  \item Isotonic regression is monotone: it re-maps probabilities without
        changing their rank order.
  \item Fitted on validation predictions and labels, applied to test
        predictions. No test leakage.
  \item Since the mapping is monotone, the threshold grid search on validation
        yields the same selected episodes before and after calibration --- the
        calibration improves interpretability of the probabilities without
        affecting the trading decision.
\end{itemize}

%%────────────────────────────────────────────────────────────────
\section{Rolling Evaluation Protocol}

\begin{table}[h]
\centering
\caption{Rolling train/validation/test folds}
\begin{tabular}{cccc}
\toprule
Fold & Train & Validation & Test \\
\midrule
1 & 2017--2019 & 2020 & 2021 \\
2 & 2018--2020 & 2021 & 2022 \\
3 & 2019--2021 & 2022 & 2023 \\
4 & 2020--2022 & 2023 & 2024 \\
5 & 2021--2023 & 2024 & 2025 \\
6 & 2022--2024 & 2025 & 2026 \\
\bottomrule
\end{tabular}
\end{table}

For each fold:
\begin{enumerate}[leftmargin=*]
  \item Scale features using \texttt{StandardScaler} fitted on training years only.
  \item Build history arrays: for each token in val/test, look back up to
        MAX\_HISTORY = 64 prior tokens from the training window.
  \item Train for up to 25 epochs with early stopping (patience = 5) on
        validation AUC. LR scheduler halves learning rate when AUC plateaus.
  \item Restore best checkpoint (highest validation AUC).
  \item Calibrate probabilities with isotonic regression on validation year.
  \item Tune threshold $\theta$ by grid search on calibrated validation P\&L.
  \item Evaluate once on test year using the validation-tuned threshold.
\end{enumerate}

%%────────────────────────────────────────────────────────────────
\section{Head Specialization}

Different heads can learn different notions of historical relevance. With 4
heads and 5 distance dimensions, different slope vectors naturally specialize:

\begin{table}[h]
\centering
\caption{Expected head specialization after training}
\begin{tabular}{lcccccc}
\toprule
Head type & $D_{\mathrm{time}}$ & $D_{\mathrm{surface}}$ &
$D_{\mathrm{short\text{-}D}}$ & $D_{\mathrm{liq}}$ & $D_{\mathrm{life}}$ \\
\midrule
Time-local   & large  & small  & small  & small  & small \\
Surface-local & small & large  & medium & small  & small \\
Short-D-local & medium & medium & large & small  & medium \\
Liquidity-local & small & small & small & large  & medium \\
\bottomrule
\end{tabular}
\end{table}

This specialization happens through gradient descent on the learned slope
parameters $\tilde{m}_h$.

%%────────────────────────────────────────────────────────────────
\section{Architecture Summary}

\begin{table}[h]
\centering
\caption{TP2-MV-ALiBi model configuration}
\begin{tabular}{ll}
\toprule
Parameter & Value \\
\midrule
Input features ($|\mathcal{F}^\star|$) & 48 \\
Hidden dimension & 64 \\
Attention heads & 4 \\
Head dimension & 16 \\
History window & 64 tokens \\
Distance dimensions & 5 \\
Strategy universe & T1, K2 \\
Main loss & BCE with class-weight \\
Auxiliary loss & Huber ($\lambda_R = 0.02$) \\
Calibration & Isotonic regression (val) \\
Optimizer & AdamW \\
Learning rate & $5 \times 10^{-4}$ \\
LR schedule & ReduceLROnPlateau (factor 0.5) \\
Epochs (max) & 25 \\
Early stopping patience & 5 \\
Batch size & 512 \\
\bottomrule
\end{tabular}
\end{table}

%%────────────────────────────────────────────────────────────────
\section{One-Line Summary}

\begin{center}
\fbox{
\begin{minipage}{0.88\textwidth}
Standard attention asks \textit{``which historical episode has similar
features?''} TP2-MV-ALiBi asks \textit{``which historical episode had a
similar option surface, similar short-D-leg risk, similar liquidity, and
similar violation lifecycle --- and also happened recently?''} The
multivariate ALiBi bias encodes all five notions of proximity directly into
the attention weights, without adding positional embeddings.
\end{minipage}
}
\end{center}

\end{document}
